%% kapa: exam=1 round=35 year=2024 period=1 form=A subject=real_estate subject_ko=부동산학원론 seq=11 number=91
%% source: Q-Net 감정평가사 기출문제 (cst00309) artlSeq=5213229 file=raw/1008/5213229/1교시 문제지 원본.pdf
%% conversion: kordoc 4.24.51 -> tools/gen1.py (auto); review=proofread-sonnet; flags=none
\begin{question}{91}
감정평가에 관한 규칙에 관한 내용으로 옳지 않은 것은?
\begin{choices}
  \item 대상물건에 대한 감정평가액은 시장가치를 기준으로 결정한다.
  \item 감정평가는 기준시점에서의 대상물건의 이용상황(불법적이거나 일시적인 이용은 제외한다) 및 공법상 제한을 받는 상태를 기준으로 한다.
  \item 감정평가는 대상물건마다 개별로 하여야 한다.
  \item 감정평가법인등이 토지를 감정평가할 때에는 수익환원법을 적용해야 한다.
  \item 하나의 대상물건이라도 가치를 달리하는 부분은 이를 구분하여 감정평가할 수 있다.
\end{choices}
\end{question}
\answer{4}
